\section{Dise\~no del \emph{front-end} anal\'ogico: topolog\'ias estudiadas y ensayadas}
\label{sec:topologias}

Esta secci\'on documenta las topolog\'ias de acondicionamiento que se estudiaron, cu\'ales se
llegaron a implementar y ensayar, y por qu\'e se descartaron las que se descartaron. Es el
n\'ucleo del trabajo de electr\'onica del proyecto.

\subsection{Planteo del problema de compensaci\'on}

Partiendo de la respuesta en aceleraci\'on del transductor, Ec.~\eqref{eq:geo-acel}:

\begin{equation}
H_a(s)=\frac{-G_0\,s}{s^{2}+2\zeta_0\omega_n s+\omega_n^{2}},
\end{equation}

el objetivo es extender la banda plana hacia frecuencias menores que $\omega_n$. Hay dos
familias de soluciones conceptualmente distintas:

\begin{enumerate}[itemsep=0.2em]
    \item \textbf{Cancelaci\'on de polos y ceros}: aplicar un filtro cuya funci\'on de
    transferencia sea la inversa de $H_a$ en la banda de inter\'es, cancelando los polos
    del transductor e imponiendo otros nuevos. Es directo pero fr\'agil: la cancelaci\'on
    exacta exige conocer $\omega_n$ y $\zeta_0$ con precisi\'on, y estos par\'ametros varían
    entre unidades, con la temperatura y con el envejecimiento. Una cancelaci\'on
    imperfecta deja un residuo que se amplifica junto con la se\~nal.
    \item \textbf{Modificaci\'on del amortiguamiento}: dejar $\omega_n$ intacta y
    \emph{aumentar} el factor de amortiguamiento de $\zeta_0$ a un valor $\zeta_1\gg\zeta_0$.
    Un sistema de segundo orden fuertemente sobreamortiguado tiene sus dos polos reales y
    muy separados, de modo que la banda entre ellos se vuelve plana. Es la estrategia
    adoptada por Ma \emph{et al.}~\cite{Ma2023} y la que se sigui\'o en este proyecto.
\end{enumerate}

\subsection{La topolog\'ia de referencia: Ma \emph{et al.} (2023)}
\label{sec:ma2023}

El trabajo de Ma \emph{et al.}~\cite{Ma2023} es la referencia principal de la cadena
anal\'ogica implementada, y merece un tratamiento detallado porque la arquitectura del
proyecto se deriva casi directamente de \'el.

\subsubsection{Principio}

El objetivo es transformar la respuesta del ge\'ofono de

\begin{equation}
H_a(s)=\frac{-G_0\,s}{s^{2}+2\zeta_0\omega_n s+\omega_n^{2}}
\qquad\text{a}\qquad
H_a'(s)=\frac{-G_0\,s}{s^{2}+2\zeta_1\omega_n s+\omega_n^{2}},
\end{equation}

es decir, aumentar el amortiguamiento manteniendo la frecuencia natural. La red de
compensaci\'on necesaria es el cociente entre ambas, y admite una descomposici\'on notable:

\begin{equation}
H_{\mathrm{comp}}(s)=\frac{H_a'(s)}{H_a(s)}
=1-\frac{2(\zeta_1-\zeta_0)\,\omega_n\,s}{s^{2}+2\zeta_1\omega_n s+\omega_n^{2}}.
\label{eq:ma-comp}
\end{equation}

La Ec.~\eqref{eq:ma-comp} es el resultado clave: la compensaci\'on requerida es
\textbf{un camino directo (pasa-todo, ganancia unitaria) menos un filtro pasa-banda de
segundo orden}. Esto se implementa con dos bloques f\'isicos:

\begin{enumerate}[itemsep=0.15em]
    \item un \textbf{filtro pasa-banda de segundo orden}, y
    \item una \textbf{etapa restadora} (sumador inversor) que sustrae la salida del
    pasa-banda de la se\~nal directa.
\end{enumerate}

\textbf{Esta descomposici\'on explica directamente la arquitectura de cuatro etapas
adoptada en este proyecto} (PGA $\to$ BP $\to$ SUM $\to$ LP, Secci\'on~\ref{sec:arq-daafe}):
las etapas BP y SUM \emph{son} la implementaci\'on f\'isica de la
Ec.~\eqref{eq:ma-comp}, con el PGA aportando la ganancia programable y el LP actuando
como \emph{anti-aliasing} antes del conversor.

\begin{figure}[ht]
\centering
\includegraphics[width=0.94\textwidth]{extension_banda.png}
\caption{Efecto ideal de elevar el amortiguamiento manteniendo $\omega_n$. La l\'inea de
\SI{5.44}{\milli\hertz} es el polo inferior del modelo con
$\zeta_1=937{,}396$; no representa una banda plana ya medida.}
\label{fig:extension-banda-main}
\end{figure}

\subsubsection{Par\'ametros del trabajo original}

\begin{center}\footnotesize
\begin{tabular}{@{}lll@{}}
\toprule
Concepto & Valor en Ma \emph{et al.}~\cite{Ma2023} & En este proyecto \\
\midrule
Ge\'ofono & JF-20DX (Sercel) & SM-24 \\
$f_n$ & \SI{10}{\hertz} & \SI{10}{\hertz} \\
Sensibilidad & \SI{2.73}{\volt\second\per\meter} & \SI{28.8}{\volt\second\per\meter} \\
Resistencia de bobina & \SI{395}{\ohm} & \SI{375}{\ohm} \\
Inductancia de bobina & \SI{50}{\milli\henry} & no especificada por el fabricante \\
Masa m\'ovil & \SI{11}{\gram} & \SI{11}{\gram} \\
$\zeta_0$ del ge\'ofono & 0{,}3 & \textbf{0{,}25} \\
$\zeta_1$ del denominador compensado & 10 & \textbf{937{,}40} nominal \\
Banda plana demostrada & \SIrange{1}{100}{\hertz} & pendiente de verificaci\'on integral \\
Topolog\'ia del pasa-banda & Sallen-Key & \textbf{MFB} (ver \S\ref{sec:mfb}) \\
\bottomrule
\end{tabular}
\end{center}

Los valores de componentes reportados para el pasa-banda son
$C_1=C_2=\SI{10}{\micro\farad}$, $R_1=\SI{78}{\ohm}$, $R_2=\SI{78}{\kilo\ohm}$
(es decir $R_2=1000R_1$) y $R_3=\SI{52}{\ohm}$ ($\approx 2R_1/3$); la etapa restadora usa
$R_4=\SI{1}{\kilo\ohm}$, $R_5=\SI{13}{\kilo\ohm}$, $R_6=\SI{100}{\kilo\ohm}$ y
$R_f=\SI{10.3}{\kilo\ohm}$, con ganancia $\times10$.

\subsubsection{Resultados reportados}

\begin{itemize}[itemsep=0.15em]
    \item Respuesta plana extendida de \SIrange{1}{100}{\hertz}, con puntos de
    \SI{-3}{\decibel} en \SI{0.5}{\hertz} y \SI{200}{\hertz}.
    \item Sensibilidad resultante $\approx\SI{194}{\milli\volt}$ por
    \si{\meter\per\second\squared} en la banda \SIrange{1}{100}{\hertz}.
    \item Mejora de relaci\'on se\~nal--ruido frente al m\'etodo de cancelaci\'on de polos y
    ceros: \SI{+1.51}{\decibel} a \SI{1}{\hertz}, \SI{+17.52}{\decibel} a
    \SI{10}{\hertz} y \SI{+21.28}{\decibel} a \SI{20}{\hertz}.
    \item Las curvas de ruido de ambos m\'etodos se mantienen por debajo del modelo de
    ruido alto de Peterson (NHNM).
\end{itemize}

La comparaci\'on con el m\'etodo de polos y ceros es justamente el argumento cuantitativo
que respalda haber elegido la familia de ``modificaci\'on del amortiguamiento'' sobre la de
``cancelaci\'on'', y es un resultado que conviene citar expl\'icitamente al justificar el
dise\~no.

\subsection{Topolog\'ias del filtro: lo que se prob\'o}
\label{sec:mfb}

\subsubsection{Filtro pasivo RC de entrada --- implementado y descartado}

La primera aproximaci\'on (17 de abril) fue un filtro pasivo RC a la entrada, con un
\emph{script} de optimizaci\'on dedicado (\texttt{FiltroPasivoEntrada.m}) para elegir los
valores de componentes. Se descart\'o en menos de \SI{24}{\hour} en favor de un FIR digital
con decimaci\'on.

\textbf{El motivo del descarte fue la tolerancia de los componentes}, y merece detalle
porque es la primera aparici\'on de un criterio que despu\'es reaparece como decisivo en
todas las decisiones de topolog\'ia del proyecto. El dise\~no optimizado quedaba con una
frecuencia de corte nominal de \SI{6.50}{\hertz}. Con tolerancias comerciales de
\SIrange{5}{10}{\percent} en resistencias y ---sobre todo--- en capacitores, esa
frecuencia var\'ia en torno a $\pm\SI{0.5}{\hertz}$ entre unidades. En un sistema de un
solo canal eso ser\'ia una molestia; en un sistema \textbf{multicanal} es un defecto
estructural, porque introduce una diferencia de fase entre canales que el algoritmo MASW
interpreta directamente como diferencia de tiempo de arribo, es decir, como velocidad de
fase err\'onea. Un filtro cuya respuesta no se puede garantizar id\'entica en todos los
canales es, en este sistema, peor que no tener filtro.

A eso se sumaron tres consideraciones secundarias, en orden de peso:

\begin{enumerate}[itemsep=0.15em]
    \item \textbf{Deriva t\'ermica.} Los valores de capacidad requeridos ---decenas de
    microfarad para cortes de unos pocos hertz--- llevan a capacitores electrol\'iticos,
    cuya capacidad var\'ia apreciablemente con la temperatura y con el envejecimiento. Un
    nodo de campo trabaja a la intemperie.
    \item \textbf{Fase no lineal} en la banda donde MASW estima velocidad de fase, con el
    mismo problema de coherencia entre canales.
    \item \textbf{El PSoC ya dispon\'ia del recurso.} El bloque de filtro digital
    programable permite realizar el mismo filtrado con coeficientes exactos,
    bit-id\'enticos en todos los nodos y de fase lineal por construcci\'on, sin costo de
    componentes adicionales.
\end{enumerate}

El reemplazo concreto fue una cadena de tres etapas FIR (\texttt{fir1(31, 0.15)}) con
decimaci\'on $\times 4$ en cada una, es decir \textbf{decimaci\'on total $\times 64$},
modelada en \texttt{FiltroDigital.slx}. La simulaci\'on que revel\'o el problema y la
sustituci\'on ocurrieron el mismo d\'ia (18 de abril).

\paragraph{Un cambio de filosof\'ia, no s\'olo de implementaci\'on.} Vale registrar que el
RC de entrada era conceptualmente un \textbf{pasa-altos por encima de la resonancia}: se
descartaba deliberadamente el contenido de baja frecuencia por considerarlo ruido. Eso es
justo lo contrario del dise\~no final, que \emph{persigue} ese contenido porque en \'el
reside la informaci\'on de las capas profundas (Secci\'on~\ref{sec:geo-park}). El descarte
del RC marca as\'i el punto en que el proyecto deja de tratar la baja frecuencia como
estorbo y pasa a tratarla como el objetivo. El RC tampoco resolv\'ia la extensi\'on de
banda del transductor, que debe actuar \emph{antes} de que la se\~nal atenuada quede bajo
el piso de ruido.

Desde entonces el reparto qued\'o fijo y se mantuvo hasta hoy: \textbf{la electr\'onica
anal\'ogica se reserva para ganancia y compensaci\'on} ---lo que necesariamente debe
ocurrir antes del conversor--- y \textbf{la selecci\'on de banda migra al dominio
digital}, donde es exacta, reproducible y reversible. La misma sensibilidad a tolerancias
que expuls\'o al RC reaparece doce d\'ias despu\'es como argumento n\'umero uno del pivote
Sallen-Key $\to$ MFB (\S\ref{sec:mfb}), y m\'as tarde como raz\'on para preferir la
compensaci\'on por amortiguamiento sobre la cancelaci\'on polo--cero
(\S\ref{sec:geo-por-que-aceleracion}). Es el hilo argumental m\'as persistente del
proyecto.

\subsubsection{Sallen-Key / VCVS --- implementado, ensayado y descartado}

Entre el 27 de abril y el 5 de mayo se estudi\'o e implement\'o la topolog\'ia VCVS
(\emph{Voltage-Controlled Voltage Source}, es decir Sallen-Key), que es la que emplea el
trabajo de referencia~\cite{Ma2023}. Se desarroll\'o el an\'alisis completo de la funci\'on
de transferencia del pasa-banda VCVS (\texttt{Compensador.m}) y se llev\'o a medici\'on.

Se descart\'o el 30 de abril tras un an\'alisis comparativo sistem\'atico de ambas
topolog\'ias sobre datos reales de captura y en simulaci\'on. El criterio de la
decisi\'on se detalla a continuaci\'on.

\subsubsection{MFB (\emph{Multiple Feedback}) --- topolog\'ia adoptada}

El 30 de abril se adopt\'o la topolog\'ia MFB, junto con la \textbf{suma en corriente}
mediante amplificador de transimpedancia. Esta \'ultima decisi\'on es coherente con la
estructura de la Ec.~\eqref{eq:ma-comp}: al ser la MFB inversora y trabajar contra una
masa virtual, la resta entre el camino directo y el pasa-banda se realiza sumando
corrientes en un nodo, sin necesidad de una etapa diferencial adicional.

\paragraph{Criterio de la decisi\'on.} Apartarse de la topolog\'ia del trabajo de
referencia~\cite{Ma2023} requiere argumento expl\'icito. Las razones que sustentaron el
cambio fueron, en orden de peso:

\begin{enumerate}[itemsep=0.3em]
    \item \textbf{Menor sensibilidad de $Q$ y $\omega_0$ a las tolerancias de los
    componentes.} Es el argumento decisivo. En la Sallen-Key pasa-banda las sensibilidades
    del factor de calidad y de la frecuencia central respecto de los valores de $R$ y $C$
    crecen con la propia $Q$, de modo que un filtro selectivo construido con componentes
    de tolerancia com\'un se aparta apreciablemente del dise\~no nominal. La MFB presenta
    sensibilidades m\'as bajas y acotadas. Dado que el proyecto se construye
    deliberadamente con \textbf{componentes de tolerancia no ajustada}, esta propiedad se
    traduce directamente en repetibilidad entre placas ---que es, adem\'as, el requisito
    de un sistema multicanal donde todos los canales deben responder igual.

    \item \textbf{Mejor atenuaci\'on en alta frecuencia.} La MFB decae de forma m\'as
    limpia fuera de la banda de paso, sin el repunte que la Sallen-Key puede presentar a
    frecuencias altas por acoplamiento directo a trav\'es de los capacitores cuando el
    amplificador pierde ganancia de lazo. Esto importa porque la etapa alimenta,
    aguas abajo, a un conversor: todo residuo de alta frecuencia se traduce en
    \emph{aliasing}.

    \item \textbf{La baja impedancia de entrada de la MFB no penaliza.} La objeci\'on
    cl\'asica contra la MFB es que su entrada carga a la etapa previa. Aqu\'i es irrelevante:
    el filtro no ve directamente al ge\'ofono sino la salida de baja impedancia de las
    etapas anteriores de la cadena. La desventaja principal de la topolog\'ia queda
    neutralizada por el orden de la arquitectura.

    \item \textbf{La inversi\'on de fase no tiene costo.} La MFB es inversora, pero en un
    sistema cuya salida se digitaliza, revertir el signo es una operaci\'on trivial en el
    dominio digital. Se trata de una desventaja nominal sin consecuencia pr\'actica.

    \item \textbf{El producto ganancia--ancho de banda no es limitante.} La MFB exige m\'as
    ganancia de lazo al amplificador que la Sallen-Key para una misma $Q$. En esta
    aplicaci\'on el margen es amplio: el GBW de los amplificadores empleados es \'ordenes de
    magnitud mayor que las frecuencias de inter\'es de las ondas Rayleigh
    (unidades a decenas de \si{\hertz}). La restricci\'on que suele decidir en contra de la
    MFB simplemente no opera en este rango.
\end{enumerate}

\begin{figure}[ht]
\centering
\includegraphics[width=0.88\textwidth]{comparacion_topologias.png}
\caption{Reconstrucci\'on trazable de \texttt{ComparacionFiltros.m}. En magnitud ambas
topolog\'ias reproducen el objetivo; la selecci\'on MFB se fundamenta en sensibilidad a
tolerancias y repetibilidad multicanal, no en una ventaja visible de respuesta nominal.}
\label{fig:comparacion-topologias-main}
\end{figure}

En s\'intesis: en la banda s\'ismica de inter\'es las desventajas de la MFB (impedancia de
entrada, inversi\'on de fase, exigencia de GBW) son inocuas por construcci\'on de la cadena
y por el rango de frecuencias, mientras que su ventaja ---robustez frente a tolerancias
de componentes--- ataca directamente el problema real del proyecto, que es construir
varios canales repetibles con componentes de bajo costo.

El antecedente m\'as cercano en la literatura reciente es el amplificador de
instrumentaci\'on con filtro diferencial de m\'ultiple realimentaci\'on (MFDFIA) propuesto
por Delis \emph{et al.}~\cite{Delis2025} para adquisici\'on de bioseñales. Aunque su
aplicaci\'on es distinta, el problema es estructuralmente el mismo que el de este
proyecto: una fuente de \textbf{se\~nal muy peque\~na con impedancia de fuente variable o
desconocida}, donde filtrar antes de amplificar es problem\'atico y filtrar despu\'es no
elimina el ruido, la interferencia de red ni los artefactos ya amplificados. La
propuesta de ese trabajo es precisamente \textbf{fusionar el filtrado con la primera
etapa de amplificaci\'on}, en lugar de tratarlos como bloques separados.

Caracter\'isticas relevantes de esa topolog\'ia:

\begin{itemize}[itemsep=0.15em]
    \item Emplea \textbf{dos amplificadores operacionales en la primera etapa}; por la
    simetr\'ia del circuito, el an\'alisis puede hacerse por separado para modo diferencial y
    modo com\'un mediante un equivalente de un solo amplificador.
    \item El rechazo de modo com\'un del conjunto queda determinado \'unicamente por el
    CMRR de la segunda etapa.
    \item En el ejemplo de aplicaci\'on reportado (adquisici\'on de EEG) se obtiene banda
    plana dentro de \SI{1}{\decibel} entre \SI{0.7}{\hertz} y \SI{25.4}{\hertz}, ganancia
    de \SI{34.1}{\decibel} y ruido referido a la entrada de \SI{70.66}{\nano\volt} rms en
    el rango \SIrange{0.1}{10}{\hertz}.
\end{itemize}

La coincidencia con este proyecto es notable en dos puntos: el PGA implementado es un
\textbf{amplificador de instrumentaci\'on de dos amplificadores
operacionales}~\cite{CypressAN60319}, y la banda de inter\'es y los niveles de se\~nal son
del mismo orden.

Este antecedente respalda, desde una aplicaci\'on independiente, las dos decisiones
centrales de la cadena implementada: amplificar antes de filtrar, y emplear m\'ultiple
realimentaci\'on en el filtro.

\subsection{Compensaci\'on: anal\'ogica frente a digital}

En paralelo con las topolog\'ias de filtro se explor\'o el problema de la compensaci\'on
propiamente dicha, con tres enfoques:

\begin{enumerate}[itemsep=0.2em]
    \item \textbf{Compensador anal\'itico} (\texttt{compensadorKai.m}, 22 de abril):
    c\'alculo directo de la red de compensaci\'on a partir de los par\'ametros del transductor.
    \item \textbf{Compensador por b\'usqueda exhaustiva} (\texttt{compensadorPazos.m}):
    barrido sobre el espacio de par\'ametros evaluando el resultado, en lugar de resolver
    anal\'iticamente. Se justifica porque la soluci\'on anal\'itica ideal suele no ser
    realizable con valores comerciales de componentes.
    \item \textbf{Cuatro motores de b\'usqueda en paralelo} (18 de mayo): cuando ninguno de
    los dos enfoques anteriores converg\'ia a una soluci\'on satisfactoria, se atac\'o el
    problema con cuatro estrategias de optimizaci\'on simult\'aneas.
\end{enumerate}

Un hallazgo relevante de esta etapa (17 de mayo) fue que \textbf{el modelo de
amplificador operacional ideal no reproduc\'ia el comportamiento medido} en el extremo
superior de la banda. La primera implementaci\'on del modelo no ideal combin\'o una
ganancia de lazo abierto $A_0=\SI{90}{\decibel}$ con el par\'ametro
$\SI{8}{\mega\hertz}$, pero lo rotul\'o err\'oneamente como polo dominante. Si
\SI{8}{\mega\hertz} es la frecuencia de ganancia unitaria de un modelo monopolo, el polo
correspondiente es $f_p=f_t/A_0\approx\SI{253}{\hertz}$, no \SI{8}{\mega\hertz}. Por
tanto, la coincidencia obtenida con esa parametrizaci\'on no se usa aqu\'i como validaci\'on
del modelo del operacional. La funci\'on de lazo abierto fue corregida en los scripts con
$f_p=f_t/A_0$; a\'un deben verificarse los par\'ametros del componente configurado y repetirse
la s\'intesis y el contraste experimental. La conclusi\'on cualitativa se mantiene ---el ancho de banda finito puede
limitar la compensaci\'on---, pero su cuantificaci\'on queda abierta.

\subsubsection{Los tres valores de amortiguamiento y la correcci\'on del prototipo}

La caracterizaci\'on de julio obliga a distinguir tres cantidades que una versi\'on previa
del borrador hab\'ia mezclado. El valor $\zeta_0=0{,}25$ es el amortiguamiento a circuito
abierto del SM-24 y, a la vez, el coeficiente que debe reproducir el \textbf{numerador}
del compensador. El denominador pasa-banda realizado con los componentes nominales
corregidos tiene $\zeta_1=937{,}396$. Finalmente, $83{,}661$ fue el amortiguamiento
residual que realizaba por error el numerador de la tanda del 20 de julio, cuando la rama
BP ten\'ia montada una resistencia incorrecta; esa tanda se archiv\'o como descartada. La
configuraci\'on corregida del 21 de julio lleva el numerador a
$\zeta=0{,}250000000000114$, coincidente con el objetivo.

La estrategia sigue siendo la de Ma \emph{et al.}~\cite{Ma2023}: aumentar el
amortiguamiento del denominador manteniendo $\omega_n$ fija. La diferencia es de grado:
Ma pasa de $0{,}3$ a $10$, mientras que el denominador nominal del prototipo alcanza
$937{,}4$. Este valor no debe confundirse con una banda plana demostrada; describe el
modelo nominal de la etapa, cuya cancelaci\'on efectiva depende de un ajuste resistivo
extremadamente fino.

La interpretaci\'on f\'isica es directa. Con $\zeta_1\gg1$ el par de polos complejos
conjugados del transductor se separa en dos polos reales situados aproximadamente en

\begin{equation}
\omega_{p1}\approx\frac{\omega_n}{2\zeta_1},
\qquad
\omega_{p2}\approx 2\zeta_1\omega_n,
\end{equation}

de modo que la banda comprendida entre ambos se vuelve plana en el modelo ideal. Con
$f_n=\SI{10.2046}{\hertz}$ y $\zeta_1=937{,}396$, la aproximaci\'on sit\'ua el polo
inferior en \SI{5.44}{\milli\hertz}. Es una capacidad \emph{te\'orica}, no una frecuencia
de corte experimental ya demostrada.

La medici\'on explica el l\'imite pr\'actico. El ajuste nominal del potenci\'ometro para
reproducir $\zeta_0=0{,}25$ es \SI{634.539}{\ohm}, mientras que la cancelaci\'on nula
ocurre en \SI{632.556}{\ohm}: apenas \SI{1.98}{\ohm} de separaci\'on sobre un recorrido
de \SI{2}{\kilo\ohm}. En $f_0$ la muesca medida queda \SI{26.33}{\decibel} menos profunda
que el modelo. Por tanto, el circuito es repetible pero la extensi\'on ideal no est\'a
verificada; el cuello de botella es la resoluci\'on del ajuste de cancelaci\'on.

Esta sensibilidad enlaza f\'isicamente la compensaci\'on con el DA-AFE. BP y SUM combinan
dos caminos de gran amplitud que deben cancelarse con precisi\'on, mientras cada etapa
aporta su propio \emph{offset}. Sin calibraci\'on, el \emph{offset} medido a la salida LP
lleg\'o a \SI{45.8}{\percent} del fondo de escala; controlar el punto de operaci\'on no es
un agregado independiente, sino lo que conserva margen din\'amico para realizar esa
resta anal\'ogica cerca del ADC.

\subsubsection{El costo en ganancia de aplanar, y por qu\'e obliga al DA-AFE}
\label{sec:costo-ganancia}

El p\'arrafo anterior enuncia el v\'inculo entre compensaci\'on y calibraci\'on como una
observaci\'on emp\'irica. En realidad es m\'as fuerte que eso: \textbf{se deduce del
\'algebra de la compensaci\'on}, y conviene demostrarlo porque es el argumento que
justifica la existencia misma del subsistema de auto-calibraci\'on. Ninguna otra parte del
documento explica por qu\'e un sistema de adquisici\'on s\'ismica necesitar\'ia un lazo de
control de \emph{offset} en cada canal.

\paragraph{La cadena compensada.} Aplicando la red de la Ec.~\eqref{eq:ma-comp} a la
respuesta del transductor, los factores $(s^2+2\zeta_0\omega_n s+\omega_n^2)$ se cancelan
y queda simplemente

\begin{equation}
H_{a,\mathrm{comp}}(s)=H_a(s)\,H_{\mathrm{comp}}(s)
=\frac{-G_0\,s}{s^{2}+2\zeta_1\omega_n s+\omega_n^{2}},
\label{eq:ha-comp}
\end{equation}

que es la misma forma de partida con $\zeta_0$ reemplazado por $\zeta_1$. Evaluando
\eqref{eq:ha-comp} en la banda media ---entre los dos polos reales, donde el t\'ermino
$2\zeta_1\omega_n s$ domina el denominador--- la sensibilidad plana resultante vale

\begin{equation}
|H_{a,\mathrm{comp}}|_{\text{plana}}\approx\frac{G_0}{2\zeta_1\omega_n},
\qquad\text{frente al pico original}\qquad
|H_a|_{\max}=\frac{G_0}{2\zeta_0\omega_n}.
\label{eq:sens-plana}
\end{equation}

\paragraph{El resultado.} Dividiendo ambas expresiones, el factor $G_0/(2\omega_n)$ se
cancela y sobrevive un cociente desnudo:

\begin{equation}
\boxed{\ \frac{|H_{a,\mathrm{comp}}|_{\text{plana}}}{|H_a|_{\max}}
=\frac{\zeta_0}{\zeta_1}\ }
\label{eq:costo-ganancia}
\end{equation}

Es decir: \textbf{aplanar la respuesta cuesta ganancia en proporci\'on directa al
amortiguamiento que se impone}. No es un efecto secundario evitable con mejor dise\~no ni
un margen que pueda recuperarse eligiendo otros componentes: es el precio algebraico de
la banda. Se est\'a canjeando ganancia por ancho de banda, exactamente como el producto
ganancia--ancho de banda de un amplificador, pero aqu\'i el intercambio es expl\'icito y
cuantificado \emph{a priori}. Num\'ericamente, para los casos de inter\'es:

\begin{center}\footnotesize
\begin{tabular}{@{}lcccc@{}}
\toprule
Caso & $\zeta_1$ & $\zeta_0/\zeta_1$ & P\'erdida [\si{\decibel}] & Polo inferior \\
\midrule
Ma \emph{et al.}~\cite{Ma2023} & $10$ & $1/40$ & $-32{,}0$ & \SI{0.5}{\hertz} \\
Tanda del 20/07 (descartada) & $83{,}661$ & $1/335$ & $-50{,}5$ & \SI{61}{\milli\hertz} \\
Nominal corregido (21/07) & $937{,}396$ & $1/3750$ & $-71{,}5$ & \SI{5.44}{\milli\hertz} \\
\bottomrule
\end{tabular}
\end{center}

\paragraph{Verificaci\'on del modelo contra el trabajo de referencia.} La
Ec.~\eqref{eq:sens-plana} puede contrastarse con datos publicados, lo que la convierte en
una relaci\'on verificada y no en una manipulaci\'on formal. Para Ma \emph{et al.}
($G_0=\SI{2.73}{\volt\second\per\meter}$, $\zeta_1=10$,
$\omega_n=2\pi\cdot\SI{10}{\hertz}=\SI{62.83}{\per\second}$) la sensibilidad plana
predicha antes de las etapas de ganancia es

\begin{equation*}
\frac{2{,}73}{2\cdot 10\cdot 62{,}83}=\SI{2.17}{\milli\volt}\ \text{por}\
\si{\meter\per\second\squared}.
\end{equation*}

El circuito publicado incorpora a continuaci\'on una etapa de ganancia $\times 10$ y una
restadora con $R_6/R_f=100/10{,}3\approx 9{,}7$, de modo que el modelo predice
$2{,}17\times 10\times 9{,}7\approx\SI{210}{\milli\volt}$ por
\si{\meter\per\second\squared}. Ma \emph{et al.} reportan \SI{194}{\milli\volt} por
\si{\meter\per\second\squared} medidos: el modelo cierra \textbf{dentro del
\SI{8}{\percent}}. De igual modo, la f\'ormula de separaci\'on de polos con $\zeta_1=10$ y
$f_n=\SI{10}{\hertz}$ predice cortes en $f_n/(2\zeta_1)=\SI{0.5}{\hertz}$ y
$2\zeta_1 f_n=\SI{200}{\hertz}$, que son \emph{exactamente} los puntos de
\SI{-3}{\decibel} que ese trabajo reporta haber medido.

\paragraph{Consecuencia 1: hay que amplificar mucho, y despu\'es a\'un m\'as.} La se\~nal
que entrega el ge\'ofono en campo ya es de pocos milivolt ---y en el extremo inferior de
la banda, de decenas de microvolt---, como se comprueba directamente en los registros
crudos (Secci\'on~\ref{sec:resultados}). Sobre esa se\~nal ya d\'ebil, la red de
compensaci\'on aplica los \SIrange{50}{70}{\decibel} de atenuaci\'on de la
Ec.~\eqref{eq:costo-ganancia}. La consecuencia pr\'actica es contraintuitiva pero
inevitable: \textbf{no es posible trabajar con ganancias moderadas}. Hay que amplificar
agresivamente \emph{antes} de la compensaci\'on, porque el ADC no ve la se\~nal
amplificada sino la se\~nal amplificada \emph{y luego atenuada} por la propia red que
aplana la banda. Toda ganancia que se ahorre en el PGA se paga en bits de resoluci\'on
perdidos en el fondo de escala del conversor. Este es el motivo real por el que la cadena
implementada dispone de PGA \emph{y} PGAout ---ganancia antes y despu\'es de la
compensaci\'on--- en lugar de una \'unica etapa de ganancia.

\paragraph{Consecuencia 2: el control de \emph{offset} deja de ser un refinamiento.} La
misma ganancia que amplifica la se\~nal amplifica el \emph{offset} de continua de cada
etapa que la precede: tensiones de \emph{offset} de entrada, corrientes de polarizaci\'on
sobre resistencias de realimentaci\'on grandes, desapareamientos del par diferencial. Con
ganancias de tres cifras, un \emph{offset} de entrada de pocos milivolt satura la etapa
siguiente antes de que llegue se\~nal alguna. La medici\'on lo confirma sin ambig\"uedad:
sin calibraci\'on, el \emph{offset} a la salida LP alcanz\'o el \SI{45.8}{\percent} del
fondo de escala ---casi la mitad del rango din\'amico del ADC consumida por una tensi\'on
que no contiene informaci\'on.

Por lo tanto, en esta cadena \textbf{el control de \emph{offset} no es un refinamiento de
calidad sino una condici\'on de existencia}: sin \'el la arquitectura de la
Ec.~\eqref{eq:ma-comp} simplemente no puede realizarse con la ganancia que la
Ec.~\eqref{eq:costo-ganancia} obliga a usar. Esto explica y justifica el subsistema
DA-AFE descrito en la Secci\'on~\ref{sec:arq-daafe}, con un VDAC/IDAC dedicado por etapa y
un lazo PI de calibraci\'on por c\'odigo de ganancia. La cadena l\'ogica completa es:

\begin{center}\footnotesize
objetivo de \SI{50}{\meter} $\Rightarrow$ banda de \SIrange{1}{2.5}{\hertz}
$\Rightarrow$ compensar por debajo de $f_n$ $\Rightarrow$ $\zeta_1\gg\zeta_0$
$\Rightarrow$ p\'erdida de \SIrange{50}{70}{\decibel} $\Rightarrow$ ganancia agresiva
$\Rightarrow$ \emph{offset} amplificado $\Rightarrow$ \textbf{DA-AFE obligatorio}
\end{center}

\paragraph{Un l\'imite superior que no fija la compensaci\'on.} Por simetr\'ia, la misma
separaci\'on de polos que lleva el corte inferior a \SI{5.44}{\milli\hertz} lleva el polo
superior a $2\zeta_1 f_n\approx\SI{19}{\kilo\hertz}$ para el valor nominal, y a
\SI{1.71}{\kilo\hertz} para $\zeta_1=83{,}661$. Ambos quedan muy por encima de la
frecuencia de Nyquist del sistema ($\SI{2604}{\hertz}/2=\SI{1302}{\hertz}$ sin
decimaci\'on, y $\SI{1302}{\hertz}/N$ con decimaci\'on). Conviene decirlo
expl\'icitamente para evitar una lectura err\'onea:
\textbf{el techo de banda de la cadena no lo fija la compensaci\'on sino el filtro
pasa-bajos \emph{anti-aliasing} y la frecuencia de muestreo}. El aporte de la
compensaci\'on es \'unicamente hacia abajo, que es donde se lo necesita.

\subsection{Filtrado de la interferencia de red}

Un problema recurrente fue la interferencia de la red el\'ectrica de \SI{50}{\hertz}, que
cae dentro de la banda de inter\'es del MASW y por lo tanto no puede simplemente filtrarse
con un pasa-bajos. Se ensayaron sucesivamente:

\begin{enumerate}[itemsep=0.2em]
    \item \textbf{FIR \emph{notch} de \SI{50}{\hertz}} (29 de abril): dise\~nado con
    \texttt{firls} (m\'inimos cuadrados con bandas ponderadas), 465 coeficientes,
    implementado en el bloque de filtro digital del PSoC. Se opt\'o por \texttt{firls}
    en lugar de \texttt{firpm} (Parks-McClellan) por el comportamiento del error en las
    bandas de transici\'on.
    \item \textbf{FIR multibanda de orden 800} (5 de mayo).
    \item \textbf{Filtro adaptativo LMS} --- ensayado y \textbf{descartado} (7 de junio) en
    favor de una soluci\'on por \textbf{m\'inimos cuadrados}, de comportamiento m\'as
    predecible. Un filtro adaptativo persigue la interferencia, pero si la se\~nal de
    inter\'es tiene contenido en la misma banda, el adaptador puede converger a cancelar
    parte de la se\~nal.
    \item \textbf{\emph{Notch} arm\'onico en el cliente web} --- implementado y luego
    \textbf{eliminado} (9 de junio) para simplificar el \emph{pipeline}; el tratamiento
    de la interferencia se traslad\'o al post-proceso, donde puede revisarse y revertirse.
\end{enumerate}

\subsubsection{Soluci\'on adoptada: b\'usqueda acotada y regresi\'on arm\'onica}

\textbf{El filtro \emph{notch} fue finalmente retirado de la cadena.} La supresi\'on de la
interferencia de red se realiza hoy en post-proceso. El algoritmo busca primero la
frecuencia fundamental en una ventana acotada alrededor de \SI{50}{\hertz}, puntuando
simult\'aneamente sus arm\'onicos. Fijada esa frecuencia, construye pares seno--coseno para
la fundamental y hasta el n\'umero de arm\'onicos elegido por el operador, estima todos los
coeficientes en una \textbf{regresi\'on lineal por m\'inimos cuadrados sobre la traza
completa} y resta la componente modelada.

El razonamiento que justifica el cambio es s\'olido y conviene explicitarlo. La
interferencia de red es una se\~nal \emph{determinista y de forma conocida}: una sinusoide
de frecuencia pr\'oxima a \SI{50}{\hertz} con amplitud y fase desconocidas, m\'as
eventualmente sus arm\'onicos. Frente a esa estructura:

\begin{itemize}[itemsep=0.2em]
    \item Un \emph{notch} de 465 coeficientes \textbf{elimina irreversiblemente} todo el
    contenido de la se\~nal en esa banda ---incluida la se\~nal s\'ismica de inter\'es, ya que
    \SI{50}{\hertz} cae dentro de la banda \'util del MASW--- e introduce un retardo de
    grupo apreciable que, en un sistema multicanal, debe ser id\'entico en todos los
    canales so pena de introducir error de velocidad de fase.
    \item La regresi\'on arm\'onica \textbf{estima} amplitud y fase de un conjunto discreto
    de sinusoides en la frecuencia de red identificada. No introduce retardo de grupo y
    evita suprimir una banda continua; adem\'as, al operar en post-proceso sobre una copia,
    es revisable y el dato crudo permanece intacto.
\end{itemize}

Se trata de reemplazar un filtrado de banda por la \textbf{estimaci\'on de una
interferencia de modelo conocido}.

\paragraph{Relaci\'on con el filtro adaptativo LMS.} La observaci\'on de que ambos
m\'etodos se parecen es correcta y conviene desarrollarla, porque la diferencia es
exactamente el motivo del cambio. En el objetivo son lo mismo: los dos estiman la
interferencia y la restan. La regresi\'on arm\'onica puede leerse, de hecho, como
\textbf{el caso l\'imite del LMS con ventana completa y modelo fijo}: si en lugar de
actualizar los pesos muestra a muestra se los resuelve de una vez sobre todo el registro,
y si en lugar de un filtro de $N$ coeficientes libres se restringe la base a un par
seno--coseno por arm\'onico, se obtiene precisamente la soluci\'on de m\'inimos cuadrados
implementada. La diferencia no est\'a en \emph{qu\'e} se calcula sino en \emph{qu\'e
libertades se le conceden al estimador}, y esas libertades son las que causaban el
problema:

\begin{center}\footnotesize
\begin{tabularx}{\textwidth}{@{}l X X@{}}
\toprule
& \textbf{LMS adaptativo} (descartado 07/06) & \textbf{M\'inimos cuadrados arm\'onicos} (adoptado) \\
\midrule
Naturaleza & Recursivo; filtro \textbf{variante en el tiempo} & Estimador \textbf{param\'etrico por lote} \\
Ventana & Deslizante, con memoria exponencial fijada por $\mu$ & \textbf{Registro completo}, sin ponderaci\'on temporal \\
R\'egimen & Atraviesa un transitorio de convergencia & Sin transitorio: no hay adaptaci\'on \\
Requisitos & Necesita una \textbf{se\~nal de referencia} de \SI{50}{\hertz} correlacionada con la interferencia & S\'olo el modelo $A\sin\omega t+B\cos\omega t$ \\
Grados de libertad & $N$ coeficientes libres del filtro & $2\times$(n\'umero de arm\'onicos) \\
Riesgo principal & Si la fuente s\'ismica entra durante la adaptaci\'on, el lazo \textbf{puede seguirla y borrarla} & El modelo es fijo: s\'olo puede restar lo que el modelo describe \\
Reversibilidad & No: modifica la traza al vuelo & S\'i: opera en post-proceso y el crudo queda intacto \\
Comportamiento multicanal & La ganancia converge distinto en cada canal $\Rightarrow$ \textbf{error de fase entre canales} & Mismos par\'ametros de modelo por canal, sin retardo de grupo \\
\bottomrule
\end{tabularx}
\end{center}

Las dos filas decisivas son las dos \'ultimas. En un sistema \textbf{multicanal} cuyo
observable es la \emph{diferencia de fase entre canales}, un filtro que converge de forma
ligeramente distinta en cada canal inyecta error directamente sobre la magnitud que se
quiere medir; y en un sistema de campo donde repetir una adquisici\'on cuesta una jornada,
la irreversibilidad de una supresi\'on mal aplicada es inaceptable. Dicho de otro modo:
\textbf{la restricci\'on del m\'etodo adoptado es su ventaja, no su limitaci\'on}. Se le
quita al estimador la libertad de perseguir la se\~nal.

La limitaci\'on residual debe quedar expl\'icita: cualquier componente s\'ismica
exactamente coherente con las sinusoides ajustadas tambi\'en se proyecta y se resta. El
m\'etodo es m\'as selectivo que un \emph{notch} ---remueve rectas espectrales en lugar de
una banda continua--- pero no garantiza preservar contenido situado exactamente en
\SI{50}{\hertz} y sus arm\'onicos.

\subsection{S\'intesis de la cadena adoptada}

\begin{center}\footnotesize
\begin{tabularx}{\textwidth}{@{}l l X@{}}
\toprule
Etapa & Topolog\'ia & Funci\'on \\
\midrule
PGA & Amp.\ de instrumentaci\'on de 2 op-amp & Ganancia programable $\times1$ a $\times50$; primera amplificaci\'on antes de cualquier filtrado \\
BP  & MFB pasa-banda & Realiza el t\'ermino sustractivo de la Ec.~\eqref{eq:ma-comp} \\
SUM & Sumador inversor (suma en corriente) & Resta el pasa-banda del camino directo: completa la compensaci\'on \\
LP  & Pasa-bajos de 2.\textsuperscript{o} orden & \emph{Anti-aliasing} previo al ADC \\
\bottomrule
\end{tabularx}
\end{center}

Cada etapa dispone de un VDAC dedicado para fijar su punto de operaci\'on, lo que habilita
la auto-calibraci\'on descrita en la Secci\'on~\ref{sec:arq-daafe}.
